\chapter{MySQL}

\section{O que é o MySQL}

O MySQL é um \textbf{SGBD relacional} de código aberto, criado em 1995 e hoje
mantido pela Oracle. É o motor que armazena e gerencia os dados; a conversa com
ele acontece em SQL.

É provavelmente o banco relacional mais usado em aplicações web: WordPress,
boa parte dos sistemas em PHP e uma infinidade de APIs rodam sobre ele. A
versão usada nesta capacitação é a \textbf{8.4 LTS}, e o motor de armazenamento
padrão é o \textbf{InnoDB}.

\begin{conceito}[Por que o InnoDB importa]
É o InnoDB que dá suporte a \textbf{chaves estrangeiras} e a \textbf{transações}.
O motor antigo, MyISAM, não suporta chave estrangeira: ele aceita a declaração e
a ignora silenciosamente. Como todo o Capítulo~\ref{cap:relacionamentos} depende
de chave estrangeira, as tabelas deste guia são criadas com
\comando{ENGINE=InnoDB} explícito.
\end{conceito}

\section{MySQL e MariaDB: a confusão mais comum}

\begin{table}[H]
\centering
\caption{Comparação entre MySQL e MariaDB}
\label{tab:mysql-mariadb}
\footnotesize
\begin{tabularx}{\textwidth}{@{}lXX@{}}
\toprule
 & \textbf{MySQL} & \textbf{MariaDB} \\
\midrule
Origem & Original, de 1995 & \emph{Fork} de 2009, feito pelos criadores do
MySQL após a compra pela Oracle \\
\addlinespace
Licença da edição livre & \emph{Community Edition}: GPLv2, gratuita &
\emph{Server}: GPLv2, gratuita \\
\addlinespace
Edições pagas & Sim (\emph{Standard} e \emph{Enterprise}, da Oracle) &
Sim (suporte e produtos da MariaDB plc) \\
\addlinespace
Compatibilidade & \multicolumn{2}{X@{}}{Altíssima nos comandos básicos: tudo
desta capacitação funciona igual nos dois} \\
\bottomrule
\end{tabularx}
\end{table}

\newpage

\begin{aviso}[Correção de um erro comum]
Dizer que ``MariaDB é livre e MySQL é comercial'' é impreciso. \textbf{Ambos têm
uma edição livre em GPLv2 e ambos têm ofertas comerciais.} A diferença que
interessa agora é outra: várias distribuições Linux (Debian, Ubuntu, Fedora,
RHEL) empacotam o \textbf{MariaDB} como padrão nos repositórios oficiais, e o
XAMPP também entrega MariaDB. Para instalar o MySQL da Oracle
especificamente, é preciso adicionar o repositório oficial do MySQL ou usar
Docker.

Para o conteúdo desta capacitação, \textbf{tanto faz}: se o seu ambiente for
MariaDB, siga normalmente.
\end{aviso}

\section{Como executar o MySQL}

Existem três caminhos. Escolha apenas um.

\subsection{Opção A: pacote tudo-em-um}

Instala servidor web, PHP, banco e phpMyAdmin de uma vez só. É o caminho mais
fácil para quem está começando.

\begin{itemize}
    \item \href{https://www.apachefriends.org/}{\textbf{XAMPP}} --- Windows,
          Linux e macOS. É o mais popular. A sigla é \textbf{X}
          (\emph{cross-platform}) + \textbf{A}pache + \textbf{M}ariaDB +
          \textbf{P}HP + \textbf{P}erl. Em versões antigas o ``M'' era MySQL;
          hoje é MariaDB.
    \item \href{https://www.wampserver.com/}{\textbf{WampServer}} --- apenas
          Windows. Aí sim: \textbf{W}indows + \textbf{A}pache + \textbf{M}ySQL
          + \textbf{P}HP.
    \item \href{https://laragon.org/}{\textbf{Laragon}} --- apenas Windows,
          alternativa mais leve e moderna.
\end{itemize}

\begin{dica}[Conflito de portas]
Se o MySQL não subir, quase sempre há outro serviço ocupando a porta 3306
(outra instalação de MySQL, ou o SQL Server). A solução é mudar a porta no
\comando{my.ini} ou \comando{my.cnf} (\comando{port=3307}) e informar essa
porta no cliente.
\end{dica}

\subsection{Opção B: instalação nativa do servidor}

\begin{itemize}
    \item \textbf{Windows:} \emph{MySQL Installer for Windows}, no site oficial.
    \item \textbf{macOS:} \comando{brew install mysql} (Homebrew) ou o
          \comando{.dmg} oficial.
    \item \textbf{Linux (Debian/Ubuntu):} \comando{sudo apt install mysql-server}.
          Atenção: em várias distribuições esse pacote resolve para MariaDB.
          Para o MySQL da Oracle, adicione antes o repositório APT oficial.
\end{itemize}

Depois de instalar, execute \comando{sudo mysql\_secure\_installation} para
definir a senha do usuário \comando{root}.

\subsection{Opção C: Docker}

\begin{aviso}[Aviso de escopo]
Docker \textbf{não} faz parte desta capacitação. Está aqui apenas para
documentar como o ambiente da aula foi montado e para estimular quem quiser
pesquisar depois. Os arquivos prontos estão na pasta \comando{infrastructure/}
do repositório.
\end{aviso}

O arquivo de orquestração é um \comando{compose.yml}, que é \textbf{diferente
de um \comando{Dockerfile}}: o \comando{Dockerfile} descreve \emph{como
construir uma imagem}; o \comando{compose.yml} descreve \emph{quais serviços
prontos subir e como conectá-los}. Como aqui só usamos imagens prontas, só
existe \comando{compose.yml}.

\begin{bashcode}[caption={Subindo o ambiente},label={lst:compose-up}]
cd infrastructure
cp .env.example .env      # preencha MYSQL_ROOT_PASSWORD
docker compose up -d      # sobe MySQL + phpMyAdmin em segundo plano
docker compose ps         # confere se estao de pe e saudaveis
\end{bashcode}

\begin{bashcode}[caption={Derrubando o ambiente},label={lst:compose-down}]
docker compose down       # para os containers, PRESERVA os dados
docker compose down -v    # para e APAGA o volume (perde tudo)
\end{bashcode}

\begin{dica}[\comando{docker compose}, sem hífen]
O comando correto é \comando{docker compose} (versão 2, um plugin do Docker).
O antigo \comando{docker-compose} com hífen é a versão 1, descontinuada
desde 2023.
\end{dica}

A área dos alunos é preparada por um script que roda automaticamente na
primeira subida do container:

\begin{sqlcode}[caption={infrastructure/init/01\_create\_databases.sql (trecho)},label={lst:init}]
-- Banco de demonstracao do instrutor: so o root tem acesso.
CREATE DATABASE IF NOT EXISTS cap CHARACTER SET utf8mb4;

-- Area dos alunos. O GRANT abaixo e sobre um PADRAO de nome, nao sobre
-- um banco que ja exista: `_` e `%` sao curingas em GRANT, por isso o
-- underscore literal precisa do escape `\_`.
-- Resultado: 'aluno' pode criar e administrar liga_ana, liga_bruno, ...
-- e nenhum outro banco.
GRANT ALL PRIVILEGES ON `liga\_%`.* TO 'aluno'@'%';
FLUSH PRIVILEGES;
\end{sqlcode}

\begin{aviso}[Scripts de \emph{init} rodam uma única vez]
Os arquivos em \comando{init/} são executados apenas na primeira subida, com o
volume de dados vazio. Se você editar um deles depois, é preciso zerar o volume
para reaplicar: \comando{docker compose down -v \&\& docker compose up -d}.
\end{aviso}
